\documentclass[10pt,a4paper]{amsart}
\usepackage[margin=19mm]{geometry}
\usepackage{amsmath,amssymb,mathtools}
\usepackage[scr=rsfs,cal=euler]{mathalfa}
\usepackage{xfrac}
\usepackage[unicode,hidelinks,bookmarks=false]{hyperref}
\newcommand{\ie}{i.e.\ }
\newcommand{\cf}{cf.\ }
\newcommand{\resp}{respectively\ }
\newcommand{\Subsection}{\subsection}
\providecommand{\nobreakdash}{\nobreak\hbox{-}}
\setlength{\emergencystretch}{2em}
\multlinegap0pt
\usepackage{tikz}
\usepackage{amssymb}
\usepackage{mathtools}
\usepackage[shortlabels]{enumitem}
\newtheorem{proposition}{Proposition}
\newcommand{\Kn}{\mathcal{K}^n}
\newcommand{\ip}[2]{\left\langle #1,#2\right\rangle}
\title{Self-dual sets up to a translation: a negative answer to Milman's question}
\author{Alex Segal}
\date{Source: arXiv:2609.12685v1 (CC BY 4.0)}
\begin{document}
\maketitle

\noindent\emph{Source and licence.} Self-dual sets up to a translation: a negative answer to Milman's question. Version arXiv:2609.12685v1 (CC BY 4.0), distributed by arXiv under the Creative Commons Attribution 4.0 International licence, http://creativecommons.org/licenses/by/4.0/, as recorded in the arXiv licence metadata for this exact version and shown on the abstract page https://arxiv.org/abs/2609.12685v1. Full licence notice: \texttt{licenses/arxiv-2609.12685-CC-BY-4.0.txt}.\par\smallskip

\noindent\textit{Editorial scope.} Counted unit: the article's proposition prop:FL\_shift (source0.tex lines 916--929). The article's complete original proof of it is delivered with it (source0.tex lines 931--1026, one \texttt{proof} environment of 96 lines). No mathematics is written, repaired or re-derived: the statement and the proof are the article's own lines, in the article's own notation, and no retained line is substituted or re-flowed.  No further source-local context is needed: every cross-reference of every retained line resolves inside the two delivered spans.  Nothing was omitted from any retained span, so the package carries no omission record.  No retained line contains a citation, so the wrapper carries no bibliography and no bibliography entry.  No retained line contains a figure environment, an image-inclusion command or drawing code, so no image is delivered and the package declares no asset dependency.  Editorial formatting only, in addition: an AMS \texttt{amsart} preamble that restates the article's own declarations and macros used by the retained lines, namely the environments \texttt{proposition} on separate counters, together with the standard packages those lines need.  The retained environments print in this document as Proposition 1 in order of appearance, whereas the article prints the same items with the numbering of its own full document.  The complete native source of the article as distributed in the arXiv e-print for this exact version is bundled unchanged as \texttt{sources/210138/source0.tex}.
\begin{proposition} \label{prop:FL_shift}
Let \(K \in \Kn\) and let \(a\in\mathbb R^n\) satisfy
\[
-a\in\operatorname{int}K^\circ.
\]Then
\[
\Phi_a(K)=(K^\circ+a)^\circ,
\]
where 
\[
\Phi_a(x):=\frac{x}{1+\langle a,x\rangle}
\]

\end{proposition}

\begin{proof}
Since \(-a\in\operatorname{int}K^\circ\), we have
\[
\langle -a,x\rangle<1
\qquad \forall x\in K,
\]
and therefore
\[
1+\langle a,x\rangle>0
\qquad \forall x\in K.
\]
Thus \(\Phi_a\) is well defined on \(K\).
Let \(x\in K\), and set $y=\Phi_a(x)$. 
For every \(z\in K^\circ\), we have $\ip{z}{x} \leq 1$, and hence
\[
\begin{aligned}
\langle z+a,y\rangle
&=
\frac{\langle z,x\rangle+\langle a,x\rangle}
     {1+\langle a,x\rangle} \\
&\leq
\frac{1+\langle a,x\rangle}
     {1+\langle a,x\rangle}
=1.
\end{aligned}
\]
It follows that $y\in(K^\circ+a)^\circ$, which implies
therefore
\(
\Phi_a(K)\subseteq(K^\circ+a)^\circ.
\)

Conversely, let
\[
y\in(K^\circ+a)^\circ.
\]
Taking \(z=0\in K^\circ\), we obtain
\(
\langle a,y\rangle\leq1.
\)
If \(\langle a,y\rangle=1\), then for every \(z\in K^\circ\),
\[
1\geq\langle z+a,y\rangle
=\langle z,y\rangle+1,
\]
so
\[
\langle z,y\rangle\leq0
\qquad \forall z\in K^\circ.
\]
Since \(0\in\operatorname{int}K^\circ\), this forces \(y=0\), contradicting
\(\ip{a}{y}=1\). Thus, it must hold that $\ip{a}{y} < 1.$

We may therefore define
\[
x:=\frac{y}{1-\langle a,y\rangle}.
\]
For every \(z\in K^\circ\),
\[
\langle z+a,y\rangle\leq1,
\]
and hence
\[
\langle z,y\rangle
\leq1-\langle a,y\rangle.
\]
Dividing by the positive number \(1-\langle a,y\rangle\), we obtain
\[
\langle z,x\rangle\leq1
\qquad \forall z \in K^\circ,
\]which in turn implies, 
\(
x\in K^{\circ\circ}=K.
\)
Finally,
\[
1+\langle a,x\rangle
=
1+\frac{\langle a,y\rangle}{1-\langle a,y\rangle}
=
\frac{1}{1-\langle a,y\rangle},
\]
and consequently
\[
\Phi_a(x)
=
\frac{x}{1+\langle a,x\rangle}
=
y.
\]
Thus
\[
(K^\circ+a)^\circ\subseteq\Phi_a(K),
\]
which proves the proposition.
\end{proof}

\end{document}
